\homework{3}{Wed 29 Oct 2025 08:57}{}

\begin{itemize}
	\item 3.1. Suppose that $X_1, X_2, \ldots$ is an i.i.d. sequence with distribution $\rho$ on some finite set. Let $\sigma_\rho^2=\operatorname{Var}\left(X_1\right)$.
a) Let $v_n$ be the sample variance of the first $n$ terms. That is,

\[
v_n=\frac{1}{n} \sum_{k=1}^n\left(X_k-\bar{X}_n\right)^2, \quad \text { where } \bar{X}_n=\frac{1}{n} \sum_{k=1}^n X_k
\]


Prove a large deviation principle for $v_n$. That is, prove that

\[
\begin{aligned}
\lim _{n \rightarrow \infty} \frac{1}{n} \log P\left(v_n>a\right) & =-\inf _{x>a} J_\rho(x), \quad \forall a>\sigma_\rho^2 \\
\text { and } \quad \lim _{n \rightarrow \infty} \frac{1}{n} \log P\left(v_n<a\right) & =-\inf _{x<a} J_\rho(x), \quad \forall 0<a<\sigma_\rho^2,
\end{aligned}
\]

where the rate function $J_\rho$ can be written as a variational expression involving the relative entropy $H(\cdot \mid \rho)$.
\\
b) In general, $J_\rho(x)$ is difficult to compute explicitly but we can compute it explicitly when $\rho=\operatorname{Bernoulli}(p)$. Compute $J_\rho(x)$ when $\rho=\operatorname{Bernoulli}(1 / 4)$.

\begin{proof}[Solution]
	By Sanov's theorem (for iid r.v. taking value in Polish spaces, see [Dembo and Zeitouni, 2010 Theorem 6.2.10), the empirical distribution  $L_n = \frac{1}{n} \sum_{i = 1}^n \delta_{X_i}$ satisfies LDP with rate $n$ and rate function $H(\cdot  | \rho)$. The sample variance is given by the following continuous map:
	\begin{align*}
		T: M_1(\mathbb{R}) &\longrightarrow \mathbb{R}^+ \\
		\nu &\longmapsto T(\nu) = \mathbb{E}_\nu (X - \mathbb{E}_ \nu X)^2
	.\end{align*}
	Then, by contraction principle, the empirical distribution of sample variance $Q_n := L_n \circ T^{-1}$ satisfies the LDP on $\mathbb{R}^+$ with rate  $n$ and rate function $J$ given by
	\[
		J(y) = \inf _{\nu \in M_1(\mathbb{R}): T \nu = y} H(\nu | \rho)
	.\] 

	For $y = \sigma_\rho^2$, we have $T \rho = y$, so $J(\sigma_\rho^2) = 0$. To obtain the tail formulas, we still need to show that $J$ is convex.

	Now consider the case $\rho = \text{Ber}(\frac{1}{4})$. In this case the support of all empirical measures in on $M_1(\{0,1\} )$. All measures $\nu$ are characterized by a single $z \in [0,1]$ by $\nu(z) = (1 - z) \delta_0 + z \delta_1$. We compute the mean, entropy and the sample variance:

	\[
		\mathbb{E}_\nu X = z
	.\] 
	\[
		\mathbb{E}_\nu (X - z)^2
		= (1 - z)^2 z + z^2 (1 - z) = z( 1- z)
	.\] 
	\[
		H(\nu | \rho) = (1 - z) \log \frac{1 - z}{\frac{3}{4}} + z \log \frac{z}{\frac{1}{4}}
	.\] 
	The rate function is given by the following constraint optimization problem
	\[
		\begin{cases}
			J(y) = \min_{z \in (0,1)} (1 - z) \log \frac{1 - z}{\frac{3}{4}} + z \log \frac{z}{\frac{1}{4}}  \\
			\text{s.t.} \quad z(1 - z) = y
		\end{cases}
	.\] 
	Using the method of Lagrange multipliers, we obtain the optimizing $z$ :
	\begin{align}
z^{*}(v)=\frac{1+\sqrt{1-4 v}}{2}
\end{align}
yielding the explicit formula
\begin{align}
J=\frac{1}{2}\left[(1-\sqrt{1-4 v}) \log \left(\frac{2(1-\sqrt{1-4 v})}{3}\right)+(1+\sqrt{1-4 v}) \log \left(\frac{1+\sqrt{1-4 v}}{8}\right)\right]
\end{align}

\end{proof}

\item 3.2. (This problem asks you to complete a step that was omitted in the proof of the contraction principle.)

Recall the setting from the proof of the contraction principle. Suppose that $T: \mathcal{X} \rightarrow \mathcal{Y}$ is a continuous function and $I$ is a rate function on $\mathcal{X}$, and let $J(y)=\inf _{x: T(x)=y} I(x)$.

Let $K_a=\{x \in \mathcal{X}: I(x) \leq a\}$ and let $L_a=\{y \in \mathcal{Y}: J(y) \leq a\}$. Prove that $L_a=T\left(K_a\right)$.
\begin{proof}[Solution]
	For each $x \in K_a$, let  $T x = y$ and note that  $J(y) \le I(x) \le a$. Hence $T K_a \subset L_a$.

	From definition of $J(y)$, for each $n \in \mathbb{N}$ we have $x_n \in K_{a + \frac{1}{n}}$ such that $T x_n = y$.
	Now, since $K_{a + \frac{1}{n}}$ is decreasing as $n \to \infty$, the whole sequence is bounded in $K_{a + 1}$, which is compact. Therefore there exists a subsequential limit $x \in \mathcal{X}$. By continuity of $I$, $x \in K_a$.

	Since $T$ is continuous on $\mathcal{X}$, 
	\[
		T x_n \to  T x = y \qquad \text{along a subsequence}
	.\] 
	This implies $y \in T K_a$, thus $L_a \subset T K_a$.
\end{proof}

\item
	3.3. A sequence of random variables $\left\{Z_n\right\}_{n \geq 1}$ is exponentially tight if for any $M<\infty$ there exists a compact set $K_M \subset \mathcal{X}$ such that

\[
\limsup _{n \rightarrow \infty} \frac{1}{n} \log P\left(Z_n \notin K_M\right) \leq-M .
\]


Prove that if $\left\{Z_n\right\}$ satisfies a weak LDP with weak rate function $I(x)$ and $\left\{Z_n\right\}$ is exponentially tight, then the weak LDP can be strengthened to a LDP with rate function $I$.
\begin{proof}[Solution]
	If $I$ does not have compact level sets in general, then there is some $a$ where the level set $K_a = \{x: I(x) \le a\} $ is closed and unbounded. Then, for any intersection with $K_M$,
	 \[
		\limsup \frac{1}{n} \log P(K_a \setminus K_M) \le  -M
	.\] 
	If we assume that the unbounded closed set $K_a$ has nonempty, unbounded interior $K_a^\circ$ (which is always true if  $I(x)$ is convex), then by weak LDP
	\[
		\liminf \frac{1}{n} \log P(K_a^\circ \setminus  K_M) \ge - I\left( K_a^\circ \setminus K_M \right) \ge - a
	.\] 
	This is gives $- M \ge  - a$, which yields contradiction for large enough $M$. Therefore, exponential tightness implies that the level sets of $I$ are in general bounded. Hence $I$ has compact level sets.

	This proof, however, breaks down for some exotic rate functions which is not convex and have infinitely many discrete points.

	Now we prove the strong LDP. Let $M > 0$. Take any closed subset $K$, and a compact subset $K_M$ as in the definition of exponential tightness. Then
	\begin{align*}
		&\limsup_{n \to \infty } \frac{1}{n} \log P_n\left( K \right) \\
		&\le \limsup_{n \to \infty } \frac{1}{n} \log \left(P_n\left( K \cap K_M \right) + P_n\left( K_M^c \right)  \right)\\
		&\leq (- I(K \cap K_M)) \vee (- M) \\
		&= - (I(K \cap K_M) \wedge M）
	.\end{align*}
	As $M \to \infty $, $K \cap K_M$ increases to $K $; by lower-semicontinuity of $I$, $I(K \cap K_M)$ converges to $- I(K)$. Therefore $- (I(K \cap K_M) \wedge M）\to  - I(K)$, as expected.
\end{proof}


		
\end{itemize}
